% Motivation and learning outcomes.
% The outcomes are the outcomes of Day 6 in the syllabus.
\begin{frame}{Why this day matters}
  \begin{itemize}
    \item Quantization makes each value smaller. It does not change the number
          of values or the number of operations.
    \item Pruning removes weights. Distillation trains a smaller model. An
          efficient design starts with fewer operations.
    \item A method that looks good on paper can give no gain on your device.
          The gain depends on the hardware and on the runtime.
    \item In the lab you prune a CNN, distill a teacher into a student,
          quantize the best models, and select one model for each board of
          this course.
  \end{itemize}
\end{frame}

\begin{frame}{Learning outcomes}
  After this day you can:
  \begin{itemize}
    \item Apply unstructured and structured magnitude pruning.
    \item Explain why sparsity does not always reduce latency.
    \item Train a small student model with knowledge distillation.
    \item Select an optimization technique from a hardware constraint.
  \end{itemize}
\end{frame}
